\section{代数叠的赋值判据: \texorpdfstring{$\Theta$}{Θ}-完备性与 S-完备性}\label{src:7.9}
\begin{epigraph}
你或许会低估定义在数学中的重要性, 但你不应该这样做.
\epigraphauthor{Claire Voisin}
\end{epigraph}
我们讨论 $\Theta$-完备性与 S-完备性的赋值判据 (见定义~\bookref{7.9.10} 和~\bookref{7.9.14}); 它们在下节的良模空间存在性定理~\bookref{7.10.1} 中起核心作用.
\subsection{从 \texorpdfstring{$\Theta=[\mathbb A^1/\mathbb G_m]$}{Θ=[A1/Gm]} 出发的映射}\label{src:7.9.1-subsection}
在 $\Spec\mathbb Z$ 上定义叠 ``Theta'' 为
\[\Theta:=[\mathbb A^1/\mathbb G_m].\]
\footnote{使用符号 $\Theta$ 是因为它的形状类似于 $\mathbb G_m$ 在复平面上的两个轨道的图像.}
在域 $\Bbbk$ 上工作时, 我们也用 $\Theta$ 表示 $\Theta_{\Bbbk}=[\mathbb A^1_{\Bbbk}/\mathbb G_{m,\Bbbk}]$. 从代数叠出发的映射 $\mathcal X\to\Theta$ 由线丛及其截面分类 (例~\bookref{4.9.25}), 而从 $\Theta$ 出发的映射 $\Theta\to\mathcal X$ 往往也有几何意义, 例如 $\mathcal X$ 是商叠, 凝聚层叠或全体曲线的叠时. 这些描述有助于理解第~\bookref{7.9.2-subsection} 节引入的 $\Theta$-完备性与 S-完备性的赋值判据.

\textbf{商叠.} 给定商叠 $[X/G]$, 单参数子群 $\lambda:\mathbb G_m\to G$ 和一点 $x\in X(\Bbbk)$, 若 $\lim_{t\to0}\lambda(t)\cdot x\in X$ 存在, 则 $\mathbb G_m$-等变延拓 $\mathbb A^1\to X$ 诱导代数叠的态射 $[\mathbb A^1/\mathbb G_m]\to[X/G]$. 下述命题断言逆命题也成立: 每个这样的映射均由单参数子群 $\lambda$ 和一点 $x\in X(\Bbbk)$ 诱导. 回忆 $C_\lambda$ 表示 $\lambda$ 的中心化子, $P_\lambda$ 是满足 $\lim_{t\to0}\lambda(t)g\lambda(t)^{-1}$ 存在的元素 $g\in G$ 组成的子群; 见第~\bookref{B.1.4-subsection} 节.
\begin{proposition}\label{src:7.9.1}
设代数闭域 $\Bbbk$ 上的光滑仿射代数群 $G$ 作用在 $\Bbbk$ 上有限型分离代数空间 $X$ 上. 则有群胚的等价
\[\operatorname{Mor}_{\Bbbk}(\Theta,[X/G])\xrightarrow{\sim}
\{(x\in X(\Bbbk),\lambda:\mathbb G_m\to G)\mid\lim_{t\to0}\lambda(t)\cdot x\in X(\Bbbk)\text{ 存在}\}.\]
态射 $(x,\lambda)\to(x',\lambda')$ 是二元组 $(g,h)$ 的同构类, 其中 $g\in P_\lambda(\Bbbk)$, $h\in G(\Bbbk)$, $x'=hgx$, $\lambda'=h\lambda h^{-1}$, 且对 $c\in C_\lambda$ 有 $(g,h)\sim(c^{-1}g,hc)$. 在此对应下, $\Theta\to[X/G]$ 将 $1$ 送到 $x$, 将 $0$ 送到 $\lim_{t\to0}\lambda(t)\cdot x$.
\end{proposition}
\begin{remark}\label{src:7.9.2}
注意 $(x,\lambda)$ 的自同构群为 $P_\lambda\cap G_x$.
\end{remark}
\begin{proof}
给定 $(x,\lambda)$, 由 $t\mapsto\lambda(t)\cdot x$ 定义的 $\mathbb G_m$-等变映射 $m_{x,\lambda}:\mathbb G_m\to X$ 延拓为交换图
\[\begin{tikzcd}[ampersand replacement=\&]\mathbb G_m\arrow[r,"{m_{x,\lambda}}"]\arrow[d,hook]\&X\\\mathbb A^1\arrow[ur,"{\widetilde m_{x,\lambda}}"']\&\end{tikzcd}\]
此延拓 $\mathbb G_m$-等变, 因而诱导商叠的态射 $f_{x,\lambda}:\Theta\to[X/G]$. 先证明这定义了函子
\begin{align}\label{src:7.9.3}
\{(x,\lambda)\mid\lim_{t\to0}\lambda(t)\cdot x\text{ 存在}\}&\longrightarrow\operatorname{Mor}_{\Bbbk}(\Theta,[X/G]),& (x,\lambda)&\longmapsto f_{x,\lambda}.
\end{align}
给定态射 $(g,h):(x,\lambda)\to(x',\lambda')$, 需要定义 $2$-态射 $f_{x,\lambda}\xrightarrow{\sim}f_{x',\lambda'}$. 因 $h$ 确定典范同构 $f_{x',\lambda'}\xrightarrow{\sim}f_{h^{-1}x',\lambda}$, 只需定义 $f_{x,\lambda}\xrightarrow{\sim}f_{h^{-1}x',\lambda}$. 因 $g\in P_\lambda(\Bbbk)$, 映射 $t\mapsto\lambda(t)g\lambda(t)^{-1}$ 延拓为 $\widetilde g:\mathbb A^1\to G$, 且 $\widetilde m_{h^{-1}x',\lambda}=\widetilde g\cdot\widetilde m_{x,\lambda}$, 因为 $h^{-1}x'=gx$. 元素 $\widetilde g$ 定义所需同构. 对 $c\in C_\lambda$, 二元组 $(g,h)$ 与 $(c^{-1}g,hc)$ 定义相同同构: 若 $c\in G_x$, 则 $(c^{-1},c)$ 定义 $f_{x,\lambda}$ 的恒等自同构. 反之, 每个同构 $f_{x,\lambda}\xrightarrow{\sim}f_{x',\lambda}$ 都由满足 $\widetilde m_{x',\lambda}=\widetilde g\cdot\widetilde m_{x,\lambda}$ 的元素 $\widetilde g\in G(\mathbb A^1)$ 诱导. 因而~\bookeqref{7.9.3} 是全忠实函子.

为证明本质满性, 设 $f:\Theta\to[X/G]$ 是态射. 在纤维图
\[\begin{tikzcd}[ampersand replacement=\&]\mathcal P\arrow[r]\arrow[d]\&X\arrow[d]\\\Theta\arrow[r,"{f}"]\&{[X/G]}\end{tikzcd}\]
中, $\mathcal P\to\Theta$ 是主 $G$-丛. 沿唯一闭点 $0\in\Theta$ 的限制 $\mathcal P|_{B\mathbb G_m}$ 对应于 $\Spec\Bbbk$ 上的 $\mathbb G_m$-等变主 $G$-丛 $P$. 选择同构 $P\cong G$ 后, $\mathcal P$ 对应于单参数子群 $\lambda:\mathbb G_m\to G$. 它在 $\mathbb A^1\times G$ 上诱导作用 $t\cdot(x,g)=(tx,g\lambda(t)^{-1})$, 从而给出主 $G$-丛
\[\mathcal P_\lambda:=[(\mathbb A^1\times G)/\mathbb G_m]\to\Theta.\]
断言存在主 $G$-丛的同构 $\alpha:\mathcal P\to\mathcal P_\lambda$. 按构造有同构 $\alpha_0:\mathcal P|_{B\mathbb G_m}\to\mathcal P_\lambda|_{B\mathbb G_m}$. 因为 $\underline{\operatorname{Isom}}_\Theta(\mathcal P,\mathcal P_\lambda)\to\Theta$ 光滑 (它是主 $G$-丛且 $G$ 光滑), 变形理论 (命题~\bookref{7.7.12}) 给出 $B\mathbb G_m\hookrightarrow\Theta$ 的幂零加厚 $\mathcal X_n$ 上相容的同构 $\alpha_n:\mathcal P|_{\mathcal X_n}\to\mathcal P_\lambda|_{\mathcal X_n}$. 凝聚 Tannaka 对偶性~\bookref{7.6.9} 与 $\Theta$ 沿 $B\mathbb G_m$ 的凝聚完备性 (定理~\bookref{7.6.13}) 将它们延拓为同构 $\alpha:\mathcal P\to\mathcal P_\lambda$. 将复合
\[\mathbb A^1\times G\to\mathcal P_\lambda\xrightarrow{\alpha^{-1}}\mathcal P\to X\]
限制到 $G$ 的单位元, 得到 $\mathbb G_m$-等变态射 $\mathbb A^1\to X$, 从而诱导映射 $f:\Theta\to[X/G]$. 令 $x$ 是 $1$ 的像, 可核实 $f_{x,\lambda}$ 与 $f$ 是 $2$-同构的.
\end{proof}
\begin{exercise}[细节]\label{src:7.9.4}
令 $R$ 为完备Noether局部 $\Bbbk$-代数, 将该命题推广为 $\operatorname{Mor}_{\Bbbk}(\Theta_R,[X/G])$ 与二元组 $(x\in X(R),\lambda:\mathbb G_m\to G)$ 之间的等价.
\end{exercise}
\begin{remark}\label{src:7.9.5}
命题~\bookref{7.9.1} 还可推广为从 $[\mathbb A^1/\mathbb G_m]$ 到 $[X/G]$ 的态射叠的描述. 即有分解
\[\underline{\operatorname{Mor}}([\mathbb A^1/\mathbb G_m],[X/G])\cong\coprod_\lambda[X_\lambda^+/P_\lambda],\]
其中 $\lambda$ 遍历单参数子群的共轭类. 此叠的代数性已在注~\bookref{7.8.13} 中讨论; 另见 \cite[定理~1.37]{HL14}.
\end{remark}
\textbf{凝聚层叠.} 对域 $\Bbbk$ 上的概形 $X$, 用 $\Coh(X)$ 表示 $X$ 上凝聚层的叠. 若 $X$ 射影 (分别地, 固有), 则由习题~\bookref{4.1.25} (分别地, 定理~\bookref{C.7.7}), $\Coh(X)$ 是代数叠.
\begin{proposition}\label{src:7.9.6}
设 $X$ 是域 $\Bbbk$ 上的概形. 对Noether $\Bbbk$-代数 $R$, $\operatorname{Mor}_{\Bbbk}(\Theta_R,\Coh(X))$ 等价于二元组 $(E,E_\bullet)$ 的群胚, 其中 $E$ 是 $X_R$ 上在 $R$ 上平坦的凝聚层,
\[E_\bullet:0\subseteq\cdots\subseteq E_{i-1}\subseteq E_i\subseteq E_{i+1}\subseteq\cdots\subseteq E\]
是 $\mathbb Z$-分次滤过, 满足 $i\ll0$ 时 $E_i=0$, $i\gg0$ 时 $E_i=E$, 且每个因子 $E_i/E_{i-1}$ 在 $R$ 上平坦. 态射 $(E,E_\bullet)\to(E',E'_\bullet)$ 是与滤过相容的凝聚层同构 $E\to E'$. 此对应将 $1$ 送到 $E$, 将 $0$ 送到相伴分次 $\operatorname{gr}E_\bullet:=\bigoplus_iE_i/E_{i-1}$. 映射经由 $\cBun(X)\subseteq\Coh(X)$ 分解当且仅当 $E$ 与每个因子 $E_i/E_{i-1}$ 都是向量丛.
\end{proposition}
\begin{proof}
读者可先阅读特殊但重要的情形 $R=\Bbbk$. 态射 $\Theta_R\to\Coh(X)$ 对应于 $X\times\Theta_R$ 上在 $\Theta_R$ 上平坦的凝聚层 $F$, 也对应于 $X\times\mathbb A_R^1$ 上在 $\mathbb A_R^1$ 上平坦, 带相容 $\mathbb G_m$-作用的凝聚层. 因为 $f:X\times\Theta_R\to X\times B\mathbb G_{m,R}$ 仿射, 可由拟凝聚层 $f_*F$ 及其 $f_*\mathcal O_{X\times\Theta_R}$-模结构恢复 $F$. $X\times B\mathbb G_{m,R}$ 上的拟凝聚层是 $X_R$ 上的 $\mathbb Z$-分次拟凝聚层 (习题~\bookref{7.1.10}), 因而 $F$ 对应于在 $R[x]$ 上平坦的 $\mathbb Z$-分次 $\mathcal O_{X_R}[x]$-模. 写成 $F=\bigoplus_{i\in\mathbb Z}E_i$, 每个 $E_i$ 是 $X_R$ 上的凝聚层, 则乘以 $x$ 诱导 $x:E_i\to E_{i+1}$. 由于 $F$ 在 $R[x]$ 上平坦, 从而无挠, 这些映射单射. $F$ 作为分次 $R[x]$-模有限生成, 所以存在有限多个次数有界的齐次生成元. 因而 $i\gg0$ 时 $E_i=E$. 另一方面, 考虑 $\mathcal O_{X_R}[x]$-子模 $E_{\geq d}:=\bigoplus_{i>d}E_i\subseteq F$. 因 $F$ Noether, 升链 $\cdots\subseteq E_{\geq d}\subseteq E_{\geq d-1}\subseteq\cdots\subseteq F$ 必终止, 因而 $i\ll0$ 时 $E_i=0$. $F$ 在 $R[x]$ 上平坦, 故商 $F/xF=\bigoplus_iE_i/E_{i-1}$ 在 $R$ 上平坦, 每个因子亦然.

反之, 给定满足条件的 $E$ 与 $E_\bullet$, 考虑分次 $\mathcal O_{X_R}[x]$-模 $F:=\bigoplus_iE_i$, 通常称为 Rees 构造. 归纳证明 $E_{\geq d}:=\bigoplus_{i\geq d}E_i$ 在 $R[x]$ 上平坦且有限生成; 因 $i\ll0$ 时 $E_i=0$, 这便证明了 $F$ 的相应性质. 当 $d\gg0$, $E_{\geq d}$ 同构于分次 $R[x]$-模 $(E\otimes_RR[x])\langle d\rangle$, 其中 $\langle d\rangle$ 表示分次移位, 所以平坦且有限生成. 对每个 $d$ 有正合列
\[0\to(E_d\otimes_RR[x])\langle d\rangle\to E_{\geq d}\to((E_{d+1}/E_d)\otimes_RR[x])\langle d+1\rangle\to0.\]
$E$ 与商 $E_{d+1}/E_d$ 的平坦性蕴含每个 $E_d$ 的平坦性. 因而左右两项在 $R[x]$ 上平坦且有限生成, 中间项亦然.
\end{proof}
\textbf{全体曲线的叠.}
\begin{proposition}\label{src:7.9.7}
设 $\mathcal M_g^{\mathrm{all}}$ 是代数闭域 $\Bbbk$ 上全体固有曲线的代数叠 (定理~\bookref{6.4.6}). 对每个 $\Bbbk$-代数 $R$, 群胚 $\operatorname{Mor}_{\Bbbk}(\Theta_R,\mathcal M_g^{\mathrm{all}})$ 的对象是 $\mathbb G_m$-等变固有曲线族 $\mathcal C\to\mathbb A_R^1$, 其中 $\mathbb G_m$ 在 $\mathbb A_R^1$ 上作通常的缩放作用; 态射是 $\mathbb G_m$-等变态射.
\end{proposition}
\begin{proof}
对 $\mathbb A_R^1\to\Theta_R$ 应用光滑下降即可.
\end{proof}
\begin{remark}\label{src:7.9.8}
其他簇的模叠也有类似描述. 文献中通常将这些 $\mathbb G_m$-等变族称为\newterm{检验构形}{test configuration}.
\end{remark}
\subsection{赋值判据: \texorpdfstring{$\Theta$}{Θ}-完备性与 S-完备性}\label{src:7.9.2-subsection}
仍记 $\Theta:=[\mathbb A^1/\mathbb G_m]$, 定义在 $\Spec\mathbb Z$ 上. 设 $R$ 是分式域为 $K$, 剩余域为 $\kappa$ 的 DVR. 定义 $\Theta_R:=\Theta\times\Spec R$, 令 $0\in\Theta_R$ 为唯一闭点. $\Theta_R$ 是权重为 $0,1$ 的商叠 $[\mathbb A^2/\mathbb G_m]$ 的局部模型: 它是良模空间 $[\mathbb A^2/\mathbb G_m]\to\Spec\Bbbk[x]$ 沿 $x\mapsto\pi$ 所给映射 $\Spec R\to\Spec\Bbbk[x]$ 的基变换, 其中 $\pi$ 是 $R$ 的一致化元. 以下笛卡尔图示意地描述 $\Theta_R$ ($x$ 是 $\mathbb A^1$ 的坐标):
\begin{align}\label{src:7.9.9}
\begin{tikzcd}[ampersand replacement=\&,column sep=small,row sep=large]
\&\Spec R\arrow[dr,hook,"{x\ne0}"]\&\&B\mathbb G_{m,R}\arrow[dl,hook',"{x=0}"']\&\\
\Spec K\arrow[ur,hook]\arrow[dr,hook]\& \&\Theta_R\&\&B\mathbb G_{m,\kappa}\arrow[ul,hook']\arrow[dl,hook]\\
\&\Theta_K\arrow[ur,hook,"{\pi\ne0}"']\&\&\Theta_\kappa\arrow[ul,hook',"{\pi=0}"]\&
\end{tikzcd}
\end{align}
左侧箭头为开浸入, 右侧箭头为闭浸入. 特别地, 到代数叠的态射 $\Theta_R\setminus0=\Spec R\cup_{\Spec K}\Theta_K\to\mathcal X$ 等价于态射 $\Spec R\to\mathcal X$ 与 $\Theta_K\to\mathcal X$, 以及它们在 $\Spec K$ 上的限制之间的同构.
\begin{definition}\label{src:7.9.10}
若对每个 DVR $R$, 每个实线箭头组成的交换图
\begin{align}\label{src:7.9.11}
\begin{tikzcd}[ampersand replacement=\&]\Theta_R\setminus0\arrow[r]\arrow[d]\&\mathcal X\\\Theta_R\arrow[ur,dashed]\&\end{tikzcd}
\end{align}
都能唯一补全, 则称Noether代数叠 $\mathcal X$ 为\newterm{$\Theta$-完备}{$\Theta$-complete} 的.\footnote{这一概念引入于 \cite{HL14}, 当时称为 ``$\Theta$-约化'' ($\Theta$-reductive).}
\end{definition}
\begin{remark}\label{src:7.9.12}
等价地, 用分类态射 $\Theta\to\mathcal X$ 的叠 $\underline{\operatorname{Mor}}(\Theta,\mathcal X)$, 在 $1$ 处的取值给出态射
\[\mathrm{ev}_1:\underline{\operatorname{Mor}}(\Theta,\mathcal X)\to\mathcal X,\qquad f\mapsto f(1).\]
$\mathcal X$ 是 $\Theta$-完备的当且仅当 $\mathrm{ev}_1$ 满足固有性的赋值判据. 若 $\mathcal X$ 在代数闭域 $\Bbbk$ 上有限型, 则 $\underline{\operatorname{Mor}}(\Theta,\mathcal X)$ 是 $\Bbbk$ 上局部有限型代数叠; 商叠情形的显式描述见注~\bookref{7.9.5}. 但此叠很少拟紧, 例如 $\mathcal X=B\mathbb G_m$ 时, 因而 $\mathrm{ev}_1$ 很少固有.
\end{remark}
对分式域为 $K$, 剩余域为 $\kappa$, 一致化元为 $\pi$ 的 DVR $R$, 定义
\begin{align}\label{src:7.9.13}
\phi_R:=[\Spec(R[s,t]/(st-\pi))/\mathbb G_m],
\end{align}
其中 $s,t$ 的 $\mathbb G_m$-权重分别为 $1,-1$.\footnote{使用符号 $\phi$ 是因为它看起来像增加了一个原点的非分离仿射直线. 此叠引入于 \cite[第~2.B 节]{Hei17}, 记作 $\overline{\mathrm{ST}}_R$, 因为它是分离性的 ``标准检验'' 概形 $\mathrm{ST}_R=\overline{\mathrm{ST}}_R\setminus0=\Spec R\cup_{\Spec K}\Spec R$ 的紧化.}
它是权重为 $1,-1$ 的商叠 $[\mathbb A^2/\mathbb G_m]$ 的局部模型, 等同于良模空间 $[\mathbb A^2/\mathbb G_m]\to\Spec\Bbbk[xy]$ 沿 $xy\mapsto\pi$ 所给映射 $\Spec R\to\Spec\Bbbk[xy]$ 的基变换. $\phi_R$ 中 $s\ne0$ 的轨迹同构于 $[\Spec(R[s,t]_s/(t-\pi/s))/\mathbb G_m]\cong[\Spec(R[s]_s)/\mathbb G_m]\cong\Spec R$, $t\ne0$ 的轨迹有类似描述. 因而有类似于~\bookeqref{7.9.9} 的笛卡尔图
\[\begin{tikzcd}[ampersand replacement=\&,column sep=small,row sep=large]
\&\Spec R\arrow[dr,hook,"{s\ne0}"]\&\&\Theta_\kappa\arrow[dl,hook',"{s=0}"']\&\\
\Spec K\arrow[ur,hook]\arrow[dr,hook]\&\&\phi_R\&\&B\mathbb G_{m,\kappa}\arrow[ul,hook']\arrow[dl,hook]\\
\&\Spec R\arrow[ur,hook,"{t\ne0}"']\&\&\Theta_\kappa\arrow[ul,hook',"{t=0}"]\&
\end{tikzcd}\]
左侧箭头是开浸入, 右侧箭头是闭浸入. 特别地, 到代数叠的态射 $\phi_R\setminus0=\Spec R\cup_{\Spec K}\Spec R\to\mathcal X$ 等价于两个态射 $\xi,\xi':\Spec R\to\mathcal X$ 及在 $\Spec K$ 上的同构 $\xi_K\simeq\xi'_K$.
\begin{definition}\label{src:7.9.14}
若对每个 DVR $R$, 每个实线箭头组成的交换图
\begin{align}\label{src:7.9.15}
\begin{tikzcd}[ampersand replacement=\&]\phi_R\setminus0\arrow[r]\arrow[d]\&\mathcal X\\\phi_R\arrow[ur,dashed]\&\end{tikzcd}
\end{align}
都能唯一补全, 则称Noether代数叠 $\mathcal X$ 为\newterm{S-完备}{S-complete} 的.\footnote{``S'' 代表 ``Seshadri'': S-完备性是一种几何性质, 令人联想到层的 S-等价关系如何蕴含模空间的分离性.}
\end{definition}
$\Theta$-完备性与 S-完备性的定义自然延伸到态射 $f:\mathcal X\to\mathcal Y$, 但我们不会使用这一推广.
\begin{exercise}\label{src:7.9.16}
设 $X$ 是 $\Bbbk$-概形, $R$ 是 $\Bbbk$ 上的 DVR.
\begin{enumerate}[(1)]
\item 证明 $X\times\phi_R$ 上的拟凝聚层 $\mathcal F$ 对应于 $X_R$ 上的 $\mathbb Z$-分次拟凝聚层 $\bigoplus_{n\in\mathbb Z}\mathcal F_n$ 及映射图
\[\begin{tikzcd}[ampersand replacement=\&]\cdots\arrow[r,bend left,"{s}"]\&\mathcal F_{n-1}\arrow[l,bend left,"{t}"]\arrow[r,bend left,"{s}"]\&\mathcal F_n\arrow[l,bend left,"{t}"]\arrow[r,bend left,"{s}"]\&\mathcal F_{n+1}\arrow[l,bend left,"{t}"]\arrow[r,bend left,"{s}"]\&\cdots\arrow[l,bend left,"{t}"]\end{tikzcd}\]
满足 $st=ts=\pi$.
\item 证明若每个 $\mathcal F_n$ 凝聚, $n\gg0$ 时 $s:\mathcal F_{n-1}\to\mathcal F_n$ 是同构, $n\ll0$ 时 $t:\mathcal F_n\to\mathcal F_{n-1}$ 是同构, 则 $\mathcal F$ 凝聚. 证明若 $s,t$ 单射, 则逆命题成立.
\item 证明 $\mathcal F$ 在 $\phi_R$ 上平坦当且仅当 $s,t$ 单射, 且诱导映射 $t:\mathcal F_{n+1}/s\mathcal F_n\to\mathcal F_n/s\mathcal F_{n-1}$ 单射.
\item 假设 $\mathcal F$ 凝聚且在 $\phi_R$ 上平坦, 如下刻画其限制:
\begin{enumerate}[(i)]
\item $X_R$ 上的 $\mathcal F|_{\{s\ne0\}}\cong\operatorname{colim}(\cdots\xrightarrow{s}\mathcal F_n\xrightarrow{s}\mathcal F_{n+1}\xrightarrow{s}\cdots)$;
\item $X_R$ 上的 $\mathcal F|_{\{t\ne0\}}\cong\operatorname{colim}(\cdots\xrightarrow{t}\mathcal F_{n+1}\xrightarrow{t}\mathcal F_n\xrightarrow{t}\cdots)$;
\item $\Theta_\kappa$ 上的 $\mathcal F|_{\{s=0\}}$ 在命题~\bookref{7.9.6} 下对应于 $\mathbb Z$-滤过 $\cdots\xrightarrow{t}\mathcal F_{n+1}/s\mathcal F_{n+2}\xrightarrow{t}\mathcal F_n/s\mathcal F_{n-1}\xrightarrow{t}\cdots$;
\item $\Theta_\kappa$ 上的 $\mathcal F|_{\{t=0\}}$ 对应于 $\mathbb Z$-滤过 $\cdots\xrightarrow{s}\mathcal F_n/t\mathcal F_{n+1}\xrightarrow{s}\mathcal F_{n+1}/t\mathcal F_{n+2}\xrightarrow{s}\cdots$.
\end{enumerate}
\end{enumerate}
\end{exercise}
\subsection{\texorpdfstring{$\Theta$}{Θ}-完备性与 S-完备性的性质}\label{src:7.9.3-subsection}
\begin{lemma}\label{src:7.9.17}
具有仿射对角态射的Noether代数叠是 $\Theta$-完备的 (分别地, S-完备的), 当且仅当每个图~\bookeqref{7.9.11} (分别地, 图~\bookeqref{7.9.15}) 都在 DVR 扩张 $R\subseteq R'$ 后存在提升. 特别地, 两种完备性均可在剩余域代数闭的完备 DVR 上检验.
\end{lemma}
\begin{proof}
先观察: 若 $\mathcal X\to\mathcal Y$ 具有仿射对角态射, $j:\mathcal U\to\mathcal T$ 是 $\mathcal Y$ 上代数叠的开浸入, 且 $j_*\mathcal O_{\mathcal U}=\mathcal O_{\mathcal T}$, 则 $\mathcal Y$-态射 $\mathcal U\to\mathcal X$ 的两个延拓 $f_1,f_2:\mathcal T\to\mathcal X$ 典范地 $2$-同构. 因为 $\underline{\operatorname{Isom}}_{\mathcal T}(f_1,f_2)\to\mathcal T$ 仿射, $2$-同构 $f_1|_{\mathcal U}\xrightarrow{\sim}f_2|_{\mathcal U}$ 诱导的截面唯一延拓为 $\mathcal T$ 上的截面.

考虑图~\bookeqref{7.9.11}, DVR 扩张 $R\subseteq R'$ 及提升 $\Theta_{R'}\to\mathcal X$. 开浸入 $j:\Theta_R\setminus0\to\Theta_R$ 满足 $j_*\mathcal O_{\Theta_R\setminus0}=\mathcal O_{\Theta_R}$; 由平坦基变换~\bookref{7.1.14}, 沿 $\Theta_{R'}\to\Theta_R$, $\Theta_{R'}\times_{\Theta_R}\Theta_{R'}\to\Theta_R$ 及三重纤维积到 $\Theta_R$ 的态射作基变换后, 同一性质仍成立. 上述观察给出两个延拓 $\Theta_{R'}\times_{\Theta_R}\Theta_{R'}\rightrightarrows\Theta_{R'}\to\mathcal X$ 之间的典范 $2$-同构, 必满足余循环条件. 由 fpqc 下降, 提升下降为 $\Theta_R\to\mathcal X$. S-完备性的论证相同.
\end{proof}
\begin{remark}\label{src:7.9.18}
当 $\mathcal X$ 在 $\Bbbk$ 上有限型时, 这些判据也可在 $\Bbbk$ 上本质有限型的 DVR 上检验; 见 \cite[第~4 节]{AHLH23}. 我们不使用这一事实.
\end{remark}
\begin{lemma}\label{src:7.9.19}
设 $f:\mathcal X\to\mathcal Y$ 是Noether代数叠的仿射态射. 若 $\mathcal Y$ 是 $\Theta$-完备的 (分别地, S-完备的), 则 $\mathcal X$ 亦然.
\end{lemma}
\begin{proof}
$\Theta_R$ 正则且 $0\in\Theta_R$ 的余维为 $2$, 因而结构层沿 $\Theta_R\setminus0\to\Theta_R$ 的推前仍是结构层. 所以有典范等价
\begin{align*}
\operatorname{Mor}_{\mathcal Y}(\Theta_R\setminus0,\mathcal X)
&\cong\operatorname{Mor}_{\mathcal O_{\mathcal Y}\text{-alg}}(f_*\mathcal O_{\mathcal X},(\Theta_R\setminus0\to\mathcal Y)_*\mathcal O_{\Theta_R\setminus0})\\
&\cong\operatorname{Mor}_{\mathcal O_{\mathcal Y}\text{-alg}}(f_*\mathcal O_{\mathcal X},(\Theta_R\to\mathcal Y)_*\mathcal O_{\Theta_R})\\
&\cong\operatorname{Mor}_{\mathcal Y}(\Theta_R,\mathcal X).
\end{align*}
S-完备性的情形完全相同.
\end{proof}
\begin{proposition}\label{src:7.9.20}
若 $G$ 是代数闭域 $\Bbbk$ 上的约化群, 则每个商叠 $[\Spec A/G]$ 都是 $\Theta$-完备且 S-完备的.
\end{proposition}
\begin{proof}
先证明 $B\mathrm{GL}_n$ 是 $\Theta$-完备的. 态射 $\Theta_R\setminus0\to\mathcal X$ 对应于 $\Theta_R\setminus0$ 上的向量丛 $E$. 叠 $\Theta_R$ 正则, $0$ 的余维为 $2$. 若 $\widetilde E$ 是延拓 $E$ 的凝聚层, 则双对偶 $\widetilde E^{\vee\vee}$ 是延拓 $E$ 的向量丛. (事实上 $E$ 沿 $\Theta_R\setminus0\hookrightarrow\Theta_R$ 的推前就是向量丛.) 这给出所需延拓 $\Theta_R\to\mathcal X$. 因 $G$ 仿射, 可选择忠实表示 $G\subseteq\mathrm{GL}_n$. 因 $G$ 约化, Matsushima 定理~\bookref{B.1.48} 蕴含 $\mathrm{GL}_n/G$ 仿射. 由笛卡尔图
\[\begin{tikzcd}[ampersand replacement=\&]\mathrm{GL}_n/G\arrow[r]\arrow[d]\&BG\arrow[d]\\\Spec\Bbbk\arrow[r]\&B\mathrm{GL}_n\end{tikzcd}\]
及光滑下降可知 $BG\to B\mathrm{GL}_n$ 仿射. 引理~\bookref{7.9.19} 蕴含 $BG$ 和 $[\Spec A/G]$ 是 $\Theta$-完备的. 命题~\bookref{7.10.30} 将证明 $BG$ 是 S-完备的, 由引理~\bookref{7.9.19}, $[\Spec A/G]$ 亦然.
\end{proof}
对环 $R$, 映射 $\Theta_R\to\Spec R$ 是良模空间, 所以由良模空间的泛性质 (定理~\bookref{7.5.11}(4)), 每个到代数空间的映射 $\Theta_A\to X$ 都经由 $\Spec R$ 分解. 当 $X$ 是 Deligne--Mumford 叠时也有类似结论.
\begin{lemma}\label{src:7.9.21}
设Noether代数叠 $\mathcal X$ 具有仿射拟有限对角态射. 若 $R$ 是完备 DVR, 则每个映射 $\Theta_R\to\mathcal X$ (分别地, $\phi_R\to\mathcal X$) 都经由 $\Theta_R\to\Spec R$ (分别地, $\phi_R\to\Spec R$) 分解.
\end{lemma}
\begin{proof}
化归到 $\mathcal X$ 仿射的情形, 此时由 $\Gamma(\Theta_R,\mathcal O_{\Theta_R})=\Gamma(\phi_R,\mathcal O_{\phi_R})=R$ 即得结论. 令 $\kappa=R/\mathfrak m$, 令 $x\in\mathcal X(\kappa)$ 为 $0\in\Theta_R$ 的像. $\mathbb G_m$ 无非平凡有限商, 所以诱导的稳定子映射 $\mathbb G_m\to G_x$ 平凡. 由命题~\bookref{5.3.4}, 可取来自仿射概形的光滑展示 $U\to\mathcal X$ 及 $x$ 的提升 $u\in U(\kappa)$. 映射 $B\mathbb G_{m,\kappa}\to\mathcal X$ 经由 $u:\Spec\kappa\to U$ 分解, 因而提升为 $B\mathbb G_{m,\kappa}\to\Spec\kappa\xrightarrow{u}U$. 令 $\mathcal T_n$ 为 $B\mathbb G_{m,\kappa}\to\Theta_R$ 的第 $n$ 个幂零加厚, 由变形理论 (命题~\bookref{7.7.12}) 可找到 $\mathcal T_n\hookrightarrow\Theta_R\to\mathcal X$ 的相容提升 $\mathcal T_n\to U$. 凝聚 Tannaka 对偶性~\bookref{7.6.9} 给出延拓 $\Theta_R\to U$. 因其经由 $\Spec R$ 分解, $\Theta_R\to\mathcal X$ 也如此.
\end{proof}
\begin{proposition}\label{src:7.9.22}
每个具有仿射拟有限对角态射的Noether代数叠 $\mathcal X$ (例如具有仿射对角态射的 Deligne--Mumford 叠) 都是 $\Theta$-完备的. 此外, $\mathcal X$ 是 S-完备的当且仅当它分离.
\end{proposition}
\begin{proof}
由引理~\bookref{7.9.17}, 可在完备 DVR $R$ 上检验两种完备性. 引理~\bookref{7.9.21} 蕴含 $\Theta$-完备性, 且蕴含 S-完备性等价于每个图
\[\begin{tikzcd}[ampersand replacement=\&]\Spec K\arrow[r]\arrow[d]\&\mathcal X\arrow[d,"{\Delta}"]\\\Spec R\arrow[r]\arrow[ur,dashed]\&\mathcal X\times\mathcal X\end{tikzcd}\]
都有提升, 这正是通常的分离性赋值判据.
\end{proof}
存在分离良模空间的必要条件是 $\Theta$-完备且 S-完备.
\begin{proposition}\label{src:7.9.23}
设 $\mathcal X$ 是代数闭域 $\Bbbk$ 上有限型, 具有仿射对角态射的代数叠. 若 $\pi:\mathcal X\to X$ 是良模空间, 则 $\mathcal X$ 是 $\Theta$-完备的. 此外, $\mathcal X$ 是 S-完备的当且仅当 $X$ 分离.
\end{proposition}
\begin{proof}
因 $X$ 是 $\Theta$-完备的, 只需证明每个实线交换图
\[\begin{tikzcd}[ampersand replacement=\&]\Theta_R\setminus0\arrow[r]\arrow[d]\&\mathcal X\arrow[d,"{\pi}"]\\\Theta_R\arrow[r]\arrow[ur,dashed]\&X\end{tikzcd}\]
存在提升. 由良模空间的局部结构~\bookref{7.7.3}, 存在包含 $\Spec R$ 的像的平展态射 $\Spec B\to X$, 满足 $\mathcal X\times_X\Spec B\cong[\Spec A/G]$, 其中 $G$ 线性约化且 $B=A^G$. $\Spec R\to X$ 在 DVR 扩张后提升到 $\Spec B$, 而完备性可在扩张后检验 (引理~\bookref{7.9.17}), 因而化归到 $[\Spec A/G]$, 即命题~\bookref{7.9.20}.

若 $X$ 分离, 则由分离性赋值判据, $\phi_R\setminus0=\Spec R\cup_{\Spec K}\Spec R\to X$ 经由 $\Spec R$ 分解, 因而 $X$ 是 S-完备的. 同样的论证说明 $\mathcal X$ 是 S-完备的. 反之, 设 $f,g:\Spec R\to X$ 满足 $f|_K=g|_K$. 适当扩张 $R$ 后, 可选择 $f|_K=g|_K$ 的提升 $\Spec K\to\mathcal X$. 因 $\mathcal X\to X$ 普遍闭 (定理~\bookref{7.5.11}(1)), 由普遍闭性的赋值判据~\bookref{4.8.7}, 进一步扩张后可取 $f,g$ 的提升 $\widetilde f,\widetilde g:\Spec R\to X$, 且 $\widetilde f|_K\cong\widetilde g|_K$. S-完备性将两者延拓为 $\phi_R\to\mathcal X$. $\phi_R\to\Spec R$ 是良模空间, 对到代数空间的映射具有泛性质, 所以 $\phi_R\to\mathcal X$ 下降为唯一的 $\Spec R\to X$, 它必同时等于 $f,g$. 分离性赋值判据即蕴含 $X$ 分离.
\end{proof}
\subsection{\texorpdfstring{$\Theta$}{Θ}-完备性与 S-完备性的例子}\label{src:7.9.4-subsection}
讨论商叠, 凝聚层叠以及全体曲线的叠的两种完备性.

\textbf{商叠.} 由命题~\bookref{7.9.1}, 映射 $\Theta\to[U/G]$ 由一点 $u\in U$ 及使 $\lim_{t\to0}\lambda(t)\cdot u$ 存在的单参数子群 $\lambda:\mathbb G_m\to G$ 分类. 据此给出商叠的 $\Theta$-完备性的几何刻画. 对 $\lambda$ 诱导的 $\mathbb G_m$-作用, 吸引子轨迹 $U_\lambda^+$ 是表示函子 $\operatorname{Mor}_{\Bbbk}^{\mathbb G_m}(\mathbb A^1,U)$ 的代数空间 (定理~\bookref{7.8.9}); 取值映射 $\mathrm{ev}_1:U_\lambda^+\to U$ 将 $f:\mathbb A^1\to U$ 送到 $f(1)$.
\begin{proposition}\label{src:7.9.24}
设 $G$ 是代数闭域 $\Bbbk$ 上的光滑线性约化群, $U$ 是 $\Bbbk$ 上有限型, 带 $G$-作用的分离代数空间. 以下条件等价:
\begin{enumerate}[(1)]
\item $[U/G]$ 是 $\Theta$-完备的;
\item 对每个来自 $\Bbbk$ 上剩余域代数闭的完备 DVR 的映射 $u:\Spec R\to U$, 以及每个单参数子群 $\lambda:\mathbb G_m\to G$, 若 $\lim_{t\to0}\lambda(t)\cdot u_K\in U(K)$ 存在, 则 $\lim_{t\to0}\lambda(t)\cdot u\in U(R)$ 也存在;
\item 对每个单参数子群 $\lambda:\mathbb G_m\to G$, 态射 $\mathrm{ev}_1:U_\lambda^+\to U$ 是闭浸入.
\end{enumerate}
\end{proposition}
\begin{proof}
因 $G$ 线性约化, $BG$ 是 $\Theta$-完备的 (命题~\bookref{7.9.20}). 因而 $[U/G]$ 的 $\Theta$-完备性等价于每个图
\begin{align}\label{src:7.9.25}
\begin{tikzcd}[ampersand replacement=\&]\Theta_R\setminus0\arrow[r]\arrow[d]\&{[U/G]}\arrow[d]\\\Theta_R\arrow[r]\arrow[ur,dashed]\&BG\end{tikzcd}
\end{align}
存在提升, 其中 $R$ 是剩余域代数闭的完备 DVR (引理~\bookref{7.9.17}). 由命题~\bookref{7.9.1} 与习题~\bookref{7.9.4}, 映射 $\Theta_R\to BG$ 对应于单参数子群 $\lambda:\mathbb G_m\to G$, 而 $\Theta_R\setminus0\to[U/G]$ 对应于使 $\lim_{t\to0}\lambda(t)\cdot u_K\in U(K)$ 存在的映射 $u:\Spec R\to U$. 换言之, 有实线箭头组成的交换图
\begin{align}\label{src:7.9.26}
\begin{tikzcd}[ampersand replacement=\&]\Spec K\arrow[r]\arrow[d]\&U_\lambda^+\arrow[d,"{\mathrm{ev}_1}"]\\\Spec R\arrow[r,"{u}"]\arrow[ur,dashed]\&U\end{tikzcd}
\end{align}
图~\bookeqref{7.9.25} 的提升对应于 $\lim_{t\to0}\lambda(t)\cdot u\in U(R)$ 的存在性, 等价于图~\bookeqref{7.9.26} 的提升. 因 $\mathrm{ev}_1$ 是有限型单态射, 它是闭浸入当且仅当固有, 当且仅当满足赋值判据的存在性部分.
\end{proof}
\begin{example}\label{src:7.9.27}
若 $U=\Spec A$ 仿射, 单参数子群 $\lambda:\mathbb G_m\to G$ 诱导分次 $A=\bigoplus_{d\in\mathbb Z}A_d$, 且 $U_\lambda^+$ 可由 $V(\sum_{d<0}A_d)$ 表示 (习题~\bookref{7.8.7}). 因而 $\mathrm{ev}_1:U_\lambda^+\hookrightarrow U$ 是闭浸入; 这再次得到 $[U/G]$ 的 $\Theta$-完备性 (命题~\bookref{7.9.20}).
\end{example}
\begin{example}\label{src:7.9.28}
考虑 $\mathbb G_m$ 在 $\mathbb P^1$ 上的作用 $t\cdot[x:y]=[tx:y]$. 取 $\lambda=\mathrm{id}$, 则 $(\mathbb P^1)_\lambda^+=\mathbb A^1\coprod\{\infty\}$, 所以 $[\mathbb P^1/\mathbb G_m]$ 非 $\Theta$-完备. 同样, 结点三次曲线 $C$ 的商叠 $[C/\mathbb G_m]$ 非 $\Theta$-完备: 若 $\mathbb P^1\to C$ 是正规化且结点上的纤维为 $\{0,\infty\}$, 则对 $\lambda=\mathrm{id}$ 有 $C_\lambda^+=\mathbb P^1\setminus\infty$. 最后, 若 $\mathbb G_m$ 通过 $t\cdot(x,y)=(tx,y)$ 作用在 $X=\mathbb A^2\setminus0$ 上, 则因 $X_\lambda^+=\{y\ne0\}$, $[X/\mathbb G_m]$ 非 $\Theta$-完备.
\end{example}
\begin{remark}\label{src:7.9.29}
态射代数叠 $\underline{\operatorname{Mor}}(\Theta,[X/G])$ 分解为不交并 $\coprod_\lambda[X_\lambda^+/P_\lambda]$, 其中 $\lambda$ 遍历单参数子群的共轭类 (注~\bookref{7.9.5}). 取值态射 $\mathrm{ev}_0:[X_\lambda^+/P_\lambda]\to[X/G]$ 由包含映射 $X_\lambda^+\to X$ 诱导. 因而 $[X/G]$ 的 $\Theta$-完备性等价于这些映射 $[X_\lambda^+/P_\lambda]\to[X/G]$ 的固有性.
\end{remark}
也可用单参数子群 $\lambda:\mathbb G_m\to G$ 及态射 $\operatorname{Mor}_{\mathbb A^1}^{\mathbb G_m}(\mathbb A^2,U\times\mathbb A^1)\to U\times U\times\mathbb A^1$ 的性质刻画 $[U/G]$ 的 S-完备性. 其中到 $U$ 的映射由沿 $x\mapsto(x,1)$ 与 $x\mapsto(1,x)$ 给出的两个映射 $\mathbb A^1\to\mathbb A^2$ 限制得到. 此判据较复杂且后文不需要, 故不在这里给出.

\textbf{凝聚层叠.} 回忆 $\Theta\to\Coh(X)$ 对应于滤过 (命题~\bookref{7.9.6}).
\begin{proposition}\label{src:7.9.30}
对代数闭域 $\Bbbk$ 上的每个射影概形 $X$, 代数叠 $\Coh(X)$ 是 $\Theta$-完备且 S-完备的.
\end{proposition}
\begin{proof}
给定 DVR $R$, 由命题~\bookref{7.9.6}, 映射 $\Theta_R\setminus0\to\Coh(X)$ 对应于 $X_R$ 上在 $R$ 上平坦的凝聚层 $E$ 和 $\mathbb Z$-分次滤过 $F_\bullet:\cdots\subseteq F_{i-1}\subseteq F_i\subseteq\cdots\subseteq E_K$, 满足 $i\gg0$ 时 $F_i=E_K$, $i\ll0$ 时 $F_i=0$, 且 $F_i/F_{i-1}$ 在 $R$ 上平坦. 将 $E$ 视为 $E_K$ 的子层, 在 $E_K$ 中定义交 $E_i:=F_i\cap E$. 因 $E_i/E_{i-1}$ 是 $F_i/F_{i-1}$ 的子层, 它无挠, 所以作为 $R$-模平坦. 滤过 $E_\bullet$ 定义延拓 $\Theta_R\to\Coh(X)$. (旁注: 这正是 Quot 概形固有性的赋值判据的论证, 见命题~\bookref{2.4.2}. 另令 $j:\Theta_R\setminus0\to\Theta_R$, $j_x:\Spec R\hookrightarrow\Theta_R$ 与 $j_\pi:\Theta_R\to\Theta_R$ 为开浸入, 令 $\mathcal E$ 是 $C\times(\Theta_R\setminus0)$ 上表示 $E$ 与 $F_\bullet$ 的并的凝聚层, 则延拓由 $(\mathrm{id}\times j)_*\mathcal E=j_{x,*}E\cap j_{\pi,*}F_\bullet=E[x^{\pm1}]\cap F_\bullet=E_\bullet$ 给出, 其中 $E[x^{\pm1}]$ 表示在每个次数都放置 $E$ 的 $\mathbb Z$-分次滤过.)

为证明 S-完备性, 设给定 $\phi_R\setminus0\to\Coh(X)$, 对应于在 $R$ 上平坦的凝聚层 $E,F$ 和同构 $\alpha:E_K\to F_K$. 回忆 $\phi_R=[\Spec(R[s,t]/(st-\pi))/\mathbb G_m]$, 有自然开浸入 $j:\phi_R\setminus0\hookrightarrow\phi_R$, $j_s,j_t:\Spec R\hookrightarrow\phi_R$ (分别满足 $s\ne0,t\ne0$), 以及 $j_{st}:\Spec K\to\phi_R$ (满足 $st\ne0$). 推前可计算为等化子
\begin{align*}
0\to(\mathrm{id}\times j)_*\mathcal E&\to(\mathrm{id}\times j_s)_*E\oplus(\mathrm{id}\times j_t)_*F\to(\mathrm{id}\times j_{st})_*F_K,\\
&&(a,b)\longmapsto a-\alpha(b).
\end{align*}
这些推前可作为 $R[s,t]/(st-\pi)$ 上的分次模计算:
\begin{align*}
(\mathrm{id}\times j_{st})_*F_K&=F_K\otimes_RR[t^{\pm1}]=\bigoplus_{n\in\mathbb Z}F_Kt^n,\\
(\mathrm{id}\times j_s)_*E&=E\otimes_RR[t^{\pm1}]=\bigoplus_{n\in\mathbb Z}Et^n,\\
(\mathrm{id}\times j_t)_*F&=F\otimes_RR[s^{\pm1}]\cong\bigoplus_{n\in\mathbb Z}(\pi^{-n}\cdot F)t^n\subseteq(\mathrm{id}\times j_{st})_*F_K,
\end{align*}
这里用了 $s=t^{-1}\pi$. 所以
\[j_*\mathcal E\cong\bigoplus_{n\in\mathbb Z}(E\cap(\pi^{-n}\cdot F))t^n\subseteq(\mathrm{id}\times j_{st})_*F_K.\]
每个 $R$-模 $E\cap(\pi^{-n}\cdot F)\subseteq E$ 都有限生成, 因 $E$ 有限生成. 升链 $\cdots\subseteq E\cap(\pi^{-n}\cdot F)\subseteq E\cap(\pi^{-n-1}\cdot F)\subseteq\cdots$ 终止于 $E$, 所以 $j_*\mathcal E$ 凝聚. 为证明它在 $\phi_R$ 上平坦, 只需检验在 $0$ 处平坦. 由平坦性的局部判据 (定理~\bookref{A.2.5}), 需证明 $\operatorname{Tor}_1^A(A/\mathfrak m,j_*\mathcal E)=0$, 其中 $A=R[s,t]/(st-\pi)$, $\mathfrak m=(s,t)$. Koszul 复形给出剩余域 $\kappa=A/\mathfrak m=R/\pi$ 的消解
\[0\to A\xrightarrow{(t,-s)}A\oplus A\xrightarrow{(s,t)}A\to\kappa\to0.\]
与 $j_*\mathcal E$ 张量得到复形
\begin{align}\label{src:7.9.31}
0\to j_*\mathcal E\xrightarrow{(t,-s)}j_*\mathcal E\oplus j_*\mathcal E\xrightarrow{(s,t)}j_*\mathcal E.
\end{align}
正合列
\[0\to\mathcal O_{\phi_R\setminus0}\xrightarrow{(t,-s)}\mathcal O_{\phi_R\setminus0}\oplus\mathcal O_{\phi_R\setminus0}\xrightarrow{(s,t)}\mathcal O_{\phi_R\setminus0}\to0\]
沿 $\mathrm{id}\times j:C\times\phi_R\setminus0\hookrightarrow C\times\phi_R$ 的推前是向量丛的左正合列, 与 $j_*\mathcal E$ 张量仍给出左正合列, 它正是~\bookeqref{7.9.31}. 所以 $\operatorname{Tor}_1^A(A/\mathfrak m,j_*\mathcal E)=0$.
\end{proof}
命题~\bookref{9.3.6} 给出 $\Coh(X)$ 的开子叠为 $\Theta$-完备且 S-完备的判据. 当 $X=C$ 为光滑曲线时, 在命题~\bookref{9.3.8} 中应用该判据证明 $\cBun_{r,d}^{\mathrm{ss}}$ 的两种完备性, 从而得出良模空间的存在性.

\textbf{全体曲线的叠.} 光滑曲线叠 $\mathcal M_g$ 与稳定曲线叠 $\overline{\mathcal M}_g$ 都是 $\Theta$-完备且 S-完备的, 因为它们是分离 Deligne--Mumford 叠. 虽然从 $\Theta$ 到全体曲线叠的映射对应于检验构形 (命题~\bookref{7.9.7}), 遗憾的是, 尚无类似上述商叠及凝聚层叠情形的已知简单判据, 能检验 $\mathcal M_g^{\mathrm{all}}$ 的给定子叠是否 $\Theta$-完备或 S-完备.
